\subsection{Security Effectiveness and Accountability}
\label{sec:evaluation-security}

We evaluate whether authenticated protocol violations are rejected and whether
audit failures yield interpretable responsibility intervals.
Table~\ref{tab:e-security-accountability} summarizes the current valid-signature
attack suite and separate historical binding and localization microbenchmarks.
These CPU experiments exercise transcript auditing independently of model inference.
Every attack query is audited; detection results therefore describe audited
requests, rather than the fraction discovered under request sampling.

\begin{table}[!t]
\centering
\caption{Security and accountability results: current valid-signature attacks
(a) and historical protocol microbenchmarks (b)--(c).}
\label{tab:e-security-accountability}
\footnotesize
\renewcommand{\arraystretch}{1.08}
\setlength{\tabcolsep}{3pt}

\begin{tabular*}{\columnwidth}{@{\hspace{3pt}}l@{\extracolsep{\fill}}ccc@{\hspace{3pt}}}
\toprule
\rowcolor{gray!12}
\multicolumn{4}{l}{\textbf{(a) Valid-signature adaptive attacks}} \\
\textbf{Scheme}
& \shortstack{\textbf{Successful}\\\textbf{episodes}}
& \shortstack{\textbf{Attack}\\\textbf{queries}}
& \shortstack{\textbf{Benign false}\\\textbf{rejections}} \\
\midrule
SignedLog & 100/100 & 100 & 0/100 \\
SignedLog+ & 0/100 & 1,600 & 0/100 \\
\textbf{\textsc{TokenProv}} & 0/100 & 1,600 & 0/100 \\
\bottomrule
\end{tabular*}

\vspace{3pt}
\begin{tabular*}{\columnwidth}{@{\hspace{3pt}}l@{\extracolsep{\fill}}cc@{\hspace{3pt}}}
\toprule
\rowcolor{gray!12}
\multicolumn{3}{l}{\textbf{(b) Cross-request splicing: accepted / tested}} \\
\textbf{Source--target pair}
& \shortstack{\textbf{Source-local}\\\textbf{checks}}
& \textbf{\textsc{TokenProv}} \\
\midrule
Same model and deployment & 500/500 & 0/500 \\
Same model, different deployment & 500/500 & 0/500 \\
Different model and deployment & 500/500 & 0/500 \\
\bottomrule
\end{tabular*}

\vspace{3pt}
\begin{tabular*}{\columnwidth}{@{\hspace{3pt}}l@{\extracolsep{\fill}}cccc@{\hspace{3pt}}}
\toprule
\rowcolor{gray!12}
\multicolumn{5}{l}{\textbf{(c) First-failure localization and responsibility}} \\
\textbf{Injected stage}
& \shortstack{\textbf{First}\\\textbf{failure}}
& \shortstack{\textbf{Relation}\\\textbf{matches}}
& \shortstack{\textbf{Interval}\\\textbf{(hops)}}
& \shortstack{\textbf{Mean audit}\\\textbf{(ms)}} \\
\midrule
Authorization & $R_{\mathrm{auth}}$ & 500/500 & $U$--$H$ (1) & 0.168 \\
Lifecycle & $R_{\mathrm{life}}$ & 500/500 & $H$--$H$ (0) & 3.476 \\
Handoff & $R_{\mathrm{hand}}$ & 500/500 & $H$--$P$ (2) & 3.430 \\
Provider & $R_{\mathrm{prov}}$ & 500/500 & $B$--$P$ (1) & 9.165 \\
Settlement & $R_{\mathrm{set}}$ & 500/500 & $H$--$B$ (1) & 14.350 \\
\bottomrule
\end{tabular*}

\vspace{2pt}
\parbox{\columnwidth}{\scriptsize
In (a), each episode permits at most 16 queries and stops at the first success.
Source-local checks in (b) use the original source context.
Each configuration in (b)--(c) uses five fixtures with 100 repetitions each.
Interval widths count hops along $U\!\rightarrow\!H\!\rightarrow\!B\!\rightarrow\!P$.}
\end{table}


\noindent\textbf{Valid-signature adaptive attacks.}
The current suite covers aggregate-bill manipulation, receipt duplication,
receipt omission, observable accepted-call omission, cross-call evidence
splicing, and unauthorized deployment substitution.
The attacker controls $H$ and $B$, and additionally controls the provider for
the splicing and deployment-substitution attacks.
It can regenerate closure and settlement records and signatures using corrupted
roles' keys; the original user authorization and retained handoff evidence
remain fixed.
Provider receipts are re-signed when the corresponding provider is corrupted.
Before each query, we check signature validity and confirm a violation against
an independent reference consisting of the original accepted calls, receipt
fields, and integer bill.
Each scheme has 100 attack episodes, with at most 16 queries per episode.
The search uses exposed first-failure feedback to select subsequent attack
families and stops at the first accepted violation.

Table~\ref{tab:e-security-accountability}(a) shows successful attacks in
100/100 episodes against \textsc{SignedLog}, which checks individual receipt
signatures. \textsc{SignedLog+} and \textsc{TokenProv} each reject all 1,600
attack queries, with no successful episode observed.
All three schemes accept their 100 benign controls.
Thus, both complete consistency checkers detect the tested violations even
when corrupted parties repair their signatures; the suite establishes no
detection-rate advantage of \textsc{TokenProv} over \textsc{SignedLog+}.
Appendix~\ref{app:e-security-accountability} provides the per-family counts
and historical fourteen-class injection results.

\noindent\textbf{Cross-request evidence binding.}
Table~\ref{tab:e-security-accountability}(b) reports a separate historical
experiment that transplants a genuine source call/receipt/artifact tuple into
a target request.
The source--target pairs use the same deployment, different deployments of
the same model, or different models and deployments.
Source-context receipt checks accept all 1,500 transplanted tuples, whereas
\textsc{TokenProv} rejects all 1,500 at the request-binding check in
$R_{\mathrm{prov}}$.
The request identifiers differ while the fixture call labels both remain
\texttt{c-0}; the result demonstrates rejection of evidence valid in its source
context but inconsistent with the audited request.
Each scenario uses five seed fixtures with 100 repetitions per fixture.

\noindent\textbf{First-failure localization.}
Table~\ref{tab:e-security-accountability}(c) covers predefined mutations in
authorization, lifecycle, handoff, provider, and settlement records.
The first failed relation matches the injected stage in 500/500 repetitions
per stage, or 2,500/2,500 overall.
The auditor maps these relations to responsibility intervals of zero to two
role-boundary hops.
Mean time to return the audit result ranges from 0.168 to 14.350\,ms across
the five stages.
These results measure stage agreement and the reported interval granularity;
an interval spanning multiple roles does not identify a unique malicious party.

\noindent\textbf{Collusion and observable evidence.}
Separate historical controls examine the effect of retaining independently
observable contradictions.
The $B+P$ splicing control, which conflicts with retained signed boundary
records, is rejected in 500/500 repetitions and returns the $B$--$P$ interval.
The $H+B$ hidden-execution control and the internally consistent $H+B+P$
control are each accepted in 500/500 repetitions when no conflicting
independent evidence is supplied.
These accepted controls identify the observability boundary and are kept
separate from the within-scope attack-success counts.
Table~\ref{tab:e-historical-collusion} describes their construction and outcomes.
